% Generated draft: requires mathematical and source review.
\documentclass[11pt]{article}
\usepackage[margin=25mm]{geometry}
\usepackage{fontspec}
\setmainfont{Latin Modern Roman}
\usepackage{amsmath,amssymb,mathtools,unicode-math}
\setmathfont{Latin Modern Math}
\usepackage{enumitem,needspace,xurl,hyperref}
\setlength{\parindent}{0pt}
\setlength{\parskip}{5pt}
\setlength{\emergencystretch}{3em}
\newcommand{\N}{\mathbb{N}}
\newcommand{\Z}{\mathbb{Z}}
\newcommand{\Q}{\mathbb{Q}}
\newcommand{\R}{\mathbb{R}}
\newcommand{\C}{\mathbb{C}}
\newcommand{\F}{\mathbb{F}}
\newcommand{\A}{\mathbb{A}}
\newcommand{\PP}{\mathbb{P}}
\newcommand{\E}{\mathbb{E}}
\newcommand{\eps}{\varepsilon}
\newcommand{\dd}{\,\mathrm{d}}
\begin{document}
\section*{UnsolvedMath ID: 30005303}
\textbf{Source problem code:} catalogue entry
\subsection{Definitions and mathematical statement}
The supplied record poses two conjectures, but does not specify the full ambient probability model or the convention for the Markov property. The statements below retain those unresolved qualifications.

Let \(G=(V,E)\) be a graph in the intended, unspecified class of graphs. Let \(\mathcal X\) be the configuration space, with coordinatewise maximum and minimum operations
\[
(x\vee y)_v=\max\{x_v,y_v\},\qquad (x\wedge y)_v=\min\{x_v,y_v\}.
\]
The coordinate domains, their orders, and the reference measure for densities are not supplied; these operations must be defined and take values in \(\mathcal X\). A probability distribution \(P\), with a specified density \(p\), is multivariate totally positive of order two, abbreviated MTP\(_2\), if
\[
p(x\vee y)p(x\wedge y)\geq p(x)p(y)
\qquad\text{for all }x,y\in\mathcal X.
\]
This is the source's pointwise requirement on the density, rather than an almost-everywhere formulation.

Write \(\mathcal M(G)\) for the source's class of distributions that are Markov with respect to \(G\), and define
\[
\mathcal M_2(G)=\{P\in\mathcal M(G):P\text{ satisfies the displayed MTP}_2\text{ condition}\}.
\]
The supplied excerpt does not specify whether the Markov property is pairwise, local, or global.

Let
\[
\mathcal A(G)=\{A\subseteq V:\{u,v\}\in E\text{ for every distinct }u,v\in A\}
\]
be the family of complete vertex subsets. The source denotes by \(\mathcal M_F(G)\) the Markov distributions whose densities factorize over complete subsets. Schematically, this means
\[
p(x)=\prod_{A\in\mathcal A(G)}\psi_A(x_A),
\]
with factor and normalization conventions requiring specification. The source denotes by \(\mathcal M_I(G)\subseteq\mathcal M_F(G)\) its subclass of Ising models, defined there by factorization using edges only:
\[
p(x)=\prod_{e\in E}\psi_e(x_e).
\]
This edge-only formula is retained as supplied. The excerpt does not explain the treatment of singleton factors, isolated vertices, or infinite products.

\textbf{Conjecture 1 (source formulation).} For every graph and probability model covered by the intended source hypotheses,
\[
\overline{\mathcal M_I(G)\cap\mathcal M_2(G)}^{\,\mathrm{pt}}
=\mathcal M_I(G)\cap\mathcal M_2(G).
\]
Here pointwise convergence is understood through the density functions. If the ambient space consists of normalized probability densities on a fixed configuration space with a common reference measure, this says that every pointwise convergent net of densities from the intersection whose limit remains an ambient probability density has its limit in the intersection. The source does not specify that ambient space, whether limits must remain normalized, or whether only sequential limits are intended.

\textbf{Conjecture 2 (source formulation).} For every graph and probability model covered by the intended source hypotheses,
\[
\mathcal M_2(G)\subseteq\mathcal M_F(G).
\]
Thus every MTP\(_2\) Markov distribution in the intended setting is conjectured to admit factorization over complete vertex subsets.

For context, the excerpt defines the extended Markov class by
\[
\mathcal M_E(G)=\overline{\mathcal M_F(G)}^{\,\mathrm{pt}}
=\overline{\mathcal M_+(G)}^{\,\mathrm{pt}},
\]
and reports
\[
\mathcal M_+(G)\subseteq\mathcal M_F(G)
\subseteq\mathcal M_E(G)\subseteq\mathcal M(G).
\]
It says that all these inclusions are strict for a generic graph, citing its reference [3], p.\ 42 ff. Neither \(\mathcal M_+(G)\), the meaning of generic, nor the bibliographic identity of [3] is supplied. It also attributes support-based sufficient conditions for a distribution in \(\mathcal M_E(G)\) to factorize to its reference [1], whose identity and hypotheses are absent. These contextual assertions are preserved as source claims, not independently verified facts.
\subsection{Short English statement}
The source asks whether the MTP\(_2\) edge-factorizing Markov models are closed under pointwise limits, and whether every MTP\(_2\) Markov distribution factorizes over complete vertex subsets. The intended graph class, state spaces, density conventions, Markov property, and ambient space for limits remain unspecified. This is a statement-only draft requiring human mathematical review.
\subsection{Sources}
\begin{itemize}
\item Local UnsolvedMath registry, \texttt{data/unsolvedmath/problems.json}: supplied record with permanent ID \(30005303\), title ``Total Positivity and Graphical Model Factorization,'' and exact source problem code \texttt{catalogue entry}. The record labels the problem open and records proposed year 2022; no proposer is named.
\item Supplied report citation: \emph{Algebraic Structures in Statistical Methodology} (2023), Oberwolfach Reports, pp.\ 3121--3170, DOI \texttt{10.4171/owr/2022/55}. Source URL: \url{https://doi.org/10.4171/owr/2022/55}. The background describes Oberwolfach Reports as workshop reports documenting open-problem discussions.
\item Supplied literature source labeled ``Lauritzen OWR report'': \url{https://doi.org/10.4171/owr/2022/55}. This supplied label does not establish authorship of the conjectures.
\item Supplied literature source labeled ``Total positivity primary article'': \url{https://par.nsf.gov/servlets/purl/10339054}. No authors, article title, or identification with the excerpt's numbered references is supplied.
\item Literature assessment dated exactly \texttt{2026-08-22}: the supplied curation labels the conjectures open, says they remain posed in the source, reports later special results for binary MTP\(_2\) models, and says no general proof or counterexample was located. This is supplied dated curation derived from the OWR material, not an independently verified literature assessment. The source itself requests independent verification of time-sensitive status and cited claims.
\item The excerpt mentions reference [1] for support-based factorization conditions and reference [3], p.\ 42 ff., for the inclusion chain, but supplies no corresponding bibliography.
\end{itemize}
\end{document}
